\documentclass{article}

\usepackage[T1]{fontenc}
\usepackage{inconsolata}
\usepackage{microtype}

\usepackage{flix}

\lstset{language=flix}

\begin{document}

\section{Algebraic Data Types and Pattern Matching}

\begin{lstlisting}
/// An algebraic data type for shapes.
enum Shape {
    case Circle(Int32),          // circle radius
    case Square(Int32),          // side length
    case Rectangle(Int32, Int32) // height and width
}

/// Computes the area of the given shape using
/// pattern matching and basic arithmetic.
def area(s: Shape): Int32 = match s {
    case Shape.Circle(r)       => 3 * (r * r)
    case Shape.Square(w)       => w * w
    case Shape.Rectangle(h, w) => h * w
}

// Computes the area of a 2 by 4.
def main(): Unit \ IO =
    println(area(Shape.Rectangle(2, 4)))
\end{lstlisting}

\section{Lists and List Processing}

\begin{lstlisting}
/// In Flix, as in many functional programming languages, 
/// lists are the bread and butter.

/// We can easily construct and append lists:
def l(): List[Int32] =
    let l1 = 1 :: 2 :: 3 :: Nil;
    let l2 = 4 :: 5 :: 6 :: Nil;
    l1 ::: l2

/// We can use pattern matching to take a list apart:
def length(l: List[a]): Int32 = match l {
  case Nil     => 0
  case _ :: xs => 1 + length(xs)
}

/// The Flix library has extensive support for lists:
def main(): Unit \ IO =
    let l1 = l();
    let l2 = List.intersperse(42, l1);
    let l3 = List.map(x -> x :: x :: Nil, l2);
    let l4 = List.flatten(l3);
    println(length(l4))
\end{lstlisting}

\section{Deriving Type Classes}
\begin{lstlisting}
/// We derive the type classes Eq, Order, and ToString
/// for the enum Month
enum Month with Eq, Order, ToString {
    case January
    case February
    case March
    case April
    case May
    case June
    case July
    case August
    case September
    case October
    case November
    case December
}

type alias Year = Int32
type alias Day = Int32

/// The Date type derives the type classes Eq and Order
enum Date(Year, Month, Day) with Eq, Order

/// We implement our own instance of ToString for Date
/// since we don't want the default `Date(1948, December, 10)`
instance ToString[Date] {
    pub def toString(x: Date): String =
        let Date.Date(y, m, d) = x;
        "${d} ${m}, ${y}"
}

/// Thanks to the Eq and Order type classes,
/// we can easily compare dates.
def earlierDate(d1: Date, d2: Date): Date = Order.min(d1, d2)

/// Thanks to the ToString type class,
/// we can easily convert dates to strings.
def printDate(d: Date): Unit \ IO =
    let message = "The date is ${d}!";
    println(message)
\end{lstlisting}

\section{Collection Literals}

\begin{lstlisting}
/// Flix has built-in literal syntax for the standard collections.
def literals(): (List[Int32], Set[Int32], Map[String, Int32], Vector[Int32], Array[Int32, Static]) =
    let xs  = List#{1, 2, 3};
    let ys  = Set#{1, 2, 3};
    let env = Map#{"min" => 0, "max" => 100};
    let zs  = Vector#{10, 20, 30};
    let arr = Array#{1, 2, 3} @ Static;
    (xs, ys, env, zs, arr)
\end{lstlisting}

\section{Restrictable Variants and Provenance}

\begin{lstlisting}
restrictable enum Color[s] {
    case Red
    case Blue
    case Green
}

/// xvar constructs extensible variants and ematch deconstructs them.
def describeColor(): String =
    let tag = xvar Color.Red;
    ematch tag {
        case Color.Red   => "red"
        case Color.Blue  => "blue"
        case Color.Green => "green"
    }

/// psolve computes a provenance model and pquery extracts one explanation.
def cycleWitness(): Vector[String] =
    let db = #{
        Edge(1, 2). Edge(2, 3). Edge(3, 1).
    };
    let rules = #{
        Path(x, y) :- Edge(x, y).
        Path(x, z) :- Path(x, y), Edge(y, z).
        Cycle(x)   :- Path(x, x).
    };
    let pm = psolve db, rules;
    pquery pm select Cycle(1) with {Edge} |>
        Vector.map(v -> ematch v { case Edge(x, y) => "${x} -> ${y}" })
\end{lstlisting}

\section{Objects and Super Calls}

\begin{lstlisting}
import java.lang.Object

/// Super calls are available inside object literals.
def mkWrappedObject(): Object \ IO =
    new Object {
        def toString(_this: Object): String \ IO =
            "wrapped:" + super.toString()

        def hashCode(_this: Object): Int32 \ IO =
            super.hashCode() + 1
    }
\end{lstlisting}

\section{Effect Handler}

\begin{lstlisting}
def main(): Unit \ { Clock, Http, Logger, IO } =
    let defaultHeaders = Map#{
        "Accept"        => List#{"application/json"},
        "Authorization" => List#{"Bearer tok123"}
    };
    run {
        let urls = List#{"/api/users", "/api/posts"};
        foreach (url <- urls) {
            match Http.get(url) {
                case Ok(resp) => println("${url} -> ${HttpResponse.status(resp)}")
                case Err(err) => println("${url} -> ${err}")
            }
        };
        match Http.get("https://notfound.flix.dev/") {
            case Ok(resp) => println("notfound -> ${HttpResponse.status(resp)}")
            case Err(err) => println("notfound -> ${err}")
        }
    } with Http.withBaseUrl("https://flix.dev")
      with Http.withDefaultHeaders(defaultHeaders)
      with Http.withRetry(Retry.linear(maxRetries = 2, delay = milliseconds(100)))
      with Http.withCircuitBreaker(failureThreshold = 3, cooldown = seconds(5))
      with Http.withSlidingWindow(maxRequests = 2, window = seconds(1))
      with Http.withLogging
\end{lstlisting}

\section{Structs and Regions}

\begin{lstlisting}
/// A mutable struct that represents a person
struct Person[r] {
    name: String,
    mut age: Int32,
    mut height: Int32
}

mod Person {
    /// Create a new `Person` struct with the specified name
    pub def giveBirth(name: String, rc: Region[r]): Person[r] \ r =
        new Person @ rc { name = name, age = 0, height = 30 }

    /// Increase the age of `person` by 1 and if the age is less than 18,
    /// increase their height by 10.
    pub def birthday(person: Person[r]): Unit \ r =
        person->age = person->age + 1;
        if (person->age < 18) {
            person->height = person->height + 10
        } else {
            ()
        }

    pub def printPerson(person: Person[r]): Unit \ {IO, r} =
        println("${person->name} is ${person->age} years old and is ${person->height} cm tall")
}

def main(): Unit \ IO = region rc {
    let joe = Person.giveBirth("Joe", rc);
    Person.birthday(joe);
    Person.birthday(joe);
    Person.birthday(joe);
    Person.printPerson(joe)
}
\end{lstlisting}

\end{document}
